\chapter{Introduction}
Legged robots can traverse environments that remain difficult or inaccessible to
conventional wheeled platforms. Their ability to establish and break contacts
with the ground allows them to negotiate obstacles, uneven surfaces, and
discontinuous terrain. This mobility, however, comes at the cost of a
substantially more complex control problem. Locomotion requires the coordinated
motion of several joints while the robot maintains balance, regulates contact
forces, and tracks the commanded velocity. These objectives must be satisfied
despite impacts, friction uncertainty, actuator limitations, and changes in the
contact configuration.\\
\
Wheeled-legged robots combine articulated limbs with driven wheels and occupy
an intermediate position between legged and wheeled systems. On regular
surfaces, rolling motion provides greater efficiency and speed than repeated
stepping. When the terrain becomes irregular, the legs can change the height
and orientation of the body, redistribute the contact forces, and step over
obstacles that cannot be crossed through rolling alone. This combination
extends the range of possible motions, but also introduces further control
challenges. The controller must account for the floating-base dynamics, the
unilateral nature of contacts, the rolling constraints of the wheels, and the
interaction between body motion and wheel traction. At higher velocities,
changes in load distribution and imperfect contact modelling can further
reduce stability \cite{racing}.\\
\
Model Predictive Control (MPC) has become one of the main approaches to locomotion
control because it expresses the task as an Optimal Control Problem (OCP). At each
control update, MPC predicts the evolution of the robot over a finite horizon
and computes a sequence of inputs that minimizes a cost function while
satisfying a model of the dynamics and a set of constraints. Only the first
input of the optimized sequence is applied. The optimization is then repeated
from the updated state. This structure allows the controller to consider future
motion rather than reacting only to the current tracking error. It also provides
a direct way to encode objectives such as velocity tracking, body-height regulation, force
distribution, friction limits, and actuator bounds.\\
\
The explicit formulation of the prediction model, objective function, and
constraints makes MPC relatively transparent. Predicted trajectories and
contact forces can be inspected, and individual weights can be adjusted
according to their physical meaning. This is particularly useful in legged
robotics, where controller failures may arise from several interacting sources.
The internal solution provides information that can help distinguish an
infeasible motion request from poor state estimation, incorrect contact
assumptions, or insufficient control authority.\\
\
This structure does not remove the difficulties associated with modelling
locomotion. A full rigid-body model describes the robot more accurately, but
its use over a sufficiently long prediction horizon may be too computationally
expensive for real-time control. Reduced-order models are therefore commonly
adopted. Examples include centroidal models, Single Rigid Body Dynamics models,
and simplified inverted-pendulum representations. These formulations retain
the quantities that dominate the global motion of the robot while omitting part
of the articulated dynamics. The reduction makes online optimization
practical, but creates a discrepancy between predicted and actual behaviour.
Joint dynamics, actuator response, structural compliance, impacts, and complex
wheel--terrain or foot--terrain interactions may only be represented
approximately.\\
\
When a reduced model is used, a Whole-Body Controller (WBC) can be introduced between
the MPC and the actuators. Its purpose is to map the motion quantities and
contact forces planned by MPC to commands that are consistent with the complete
kinematics and dynamics of the robot. A WBC can be formulated as a Quadratic
Programming (QP) problem in which generalized accelerations, contact forces, or joint torques
are selected to track the reduced-order plan. At the same time, the QP can
account for additional requirements that are absent from the simplified MPC
model, including the floating-base equations of motion, friction-cone
conditions, joint limits, contact constraints, swing-foot tracking,
body-orientation regulation, and posture regularization. The resulting
hierarchy separates long-horizon motion planning from instantaneous whole-body
control. It does not eliminate model mismatch, but it allows several physical
constraints to be enforced without introducing the complete multibody system
into every node of the MPC horizon.\\
\
The performance of the resulting MPC-WBC architecture still depends on the
validity of its models and contact assumptions. Errors in mass and inertia
parameters, terrain height, friction, actuation delay, or contact timing can
produce persistent tracking errors and, in more severe cases, instability.
Some effects can be addressed by increasing model fidelity or adding
constraints and cost terms. This solution is not always practical. A more
detailed formulation increases the computational burden, and some phenomena
are difficult to describe through a compact differentiable model. The
controller therefore faces a persistent compromise between physical accuracy
and the ability to solve the optimization problem at the required frequency.\\
\
Most established MPC solvers used in robotics are designed primarily for CPU
execution. Although these implementations can achieve the update rates required
to control a single robot, their computational cost becomes a bottleneck when
the controller must be evaluated across thousands of parallel simulation
environments. This requirement is particularly relevant when MPC is included
directly in a Reinforcement Learning training loop. Recent work has therefore
focused on reformulating optimal-control algorithms to exploit GPU parallelism,
allowing multiple MPC problems to be solved concurrently alongside the learning
process. MPX follows this direction by providing a JAX-based, GPU-accelerated
MPC framework for legged robots \cite{mpx}. Its solver exploits temporal and
state-space parallelism through associative scan operations, enabling the
model-based controller to remain active during large-scale policy training
rather than being used only to generate offline trajectories or reference
motions. Its solver exploits temporal and
state-space parallelism through associative scan operations, reducing the
dependence of computation time on the prediction-horizon length and on the
dimensions of the state and control spaces. Besides supporting real-time
control, this formulation makes it possible to evaluate many MPC instances
concurrently during learning. Model-based control can therefore remain inside
the simulation loop rather than being used only to generate an offline dataset
or a reference trajectory.\\
\
Reinforcement Learning (RL) offers a different approach to the same control problem.
Rather than requiring a complete analytical description of the dynamics, an RL
policy learns a mapping from observations to actions through interaction with
an environment. The policy is optimized to maximize the expected cumulative
reward, which can include velocity tracking, stability, energy consumption,
action smoothness, and other task-related objectives. Training under variations
in the robot model, terrain, observations, and external disturbances can
produce behaviours that are less dependent on a single nominal model. RL can
also account implicitly for nonlinear contact effects and objectives that would
be cumbersome to include in an online optimization problem.\\
\
These advantages are accompanied by important limitations. A learned policy
does not normally enforce physical constraints in the explicit form used by
MPC. Its behaviour depends strongly on the training distribution, reward
definition, action representation, and simulation model. Good performance
inside the training domain does not guarantee predictable behaviour in states
that were rarely or never encountered. Reward design can also require
considerable tuning, since a policy may exploit unintended aspects of the
simulator or trade one objective against another in an undesirable way.
Finally, the internal reasoning of a neural policy is difficult to inspect.
Diagnosing a failure is therefore less direct than examining the state,
constraints, and cost terms of a model-based optimizer.\\
\
MPC and RL have complementary strengths. MPC supplies a physically grounded
nominal behaviour, explicit predictions, and structured constraint handling.
RL provides a mechanism for learning corrections from experience when the
nominal model is incomplete. This complementarity has motivated several
combined architectures. A policy may adjust parameters of the optimization
problem, generate reference trajectories or contact schedules, imitate
trajectories produced by an MPC expert, or act in parallel with the model-based
controller \cite{adaptive_mpc,non_gaited}. Each choice
determines how authority is divided between the analytical controller and the
learned component.\\
\
This thesis focuses on residual learning, in which the model-based controller
remains active and the policy produces a correction to its nominal command. The
main locomotion task is therefore not learned from an unstructured initial
policy. MPC and Whole-Body Control provide the nominal motion and the
corresponding actuator torques, while RL learns only the additional contribution
required to improve the closed-loop response. The final command can be
expressed in the general form

$$
\boldsymbol{\tau}_{\mathrm{cmd}}
=
\boldsymbol{\tau}_{\mathrm{nom}}
+
\boldsymbol{\tau}_{\mathrm{res}},
$$

where \(\boldsymbol{\tau}_{\mathrm{nom}}\) is generated by the model-based
control branch and \(\boldsymbol{\tau}_{\mathrm{res}}\) is obtained from the
policy output through a low-level feedback law. The nominal controller provides
a strong control prior and keeps the learning process close to a physically
meaningful locomotion strategy. The policy is then free to compensate for
systematic tracking errors, simplified dynamics, uncertain contacts, and other
effects that are not captured adequately by the nominal formulation. Similar
torque-level residual architectures have shown that an MPC prior can improve
training efficiency and extend the range of trackable locomotion commands
\cite{residual_mpc}.\\
\
The residual structure also provides a natural experimental baseline. When the
learned contribution is bypassed, the robot is controlled by the original architecture. 
When it is enabled, the nominal controller and
experimental conditions remain unchanged, while the policy adds its correction.
The effect of learning can therefore be evaluated without comparing two
entirely different locomotion pipelines. This distinction must be implemented
carefully: setting the policy action to zero does not necessarily remove every
torque produced by a residual feedback controller. A genuine baseline requires
the residual branch to be excluded from the actuator command.\\
\
The work presented in this thesis studies this approach on two robotic systems
in the MuJoCo physics engine \cite{mujoco}. MuJoCo provides
rigid-body dynamics, contact simulation, and an interface suitable for parallel
learning and control experiments. Both the nominal controllers and the learned
policies are evaluated within the same simulated dynamics, allowing their
behaviour to be compared under consistent initial conditions, velocity
commands, disturbances, and model variations.\\
\
The first platform is TITA, a wheeled-legged robot with two articulated legs
and two actuated wheels. Its nominal MPC is based on the Double Force
Constrained Inverted Pendulum model (DFCIP). This reduced-order
representation describes the quantities that dominate the motion of the robot,
including its centre-of-mass position and velocity, heading, forward velocity,
yaw rate, and wheel contact geometry. From the measured robot state and the
commanded forward and angular velocities, the MPC plans the desired forward,
vertical, and angular accelerations together with the ground reaction forces at
the two wheels.\\
\
These accelerations and ground reaction forces constitute the reduced-order
references supplied by the MPC to the Whole-Body Controller. The WBC formulates
an instantaneous Quadratic Program over the generalized accelerations and wheel
contact forces. Its objective includes tracking tasks for the centre of mass,
the left and right wheels, the base orientation, and the joint posture. The
optimization also accounts for the floating-base equations of motion, friction
constraints, and predicted joint limits. Once the QP has been solved, inverse
dynamics maps the optimized generalized accelerations and contact forces to the
nominal torques of the eight actuators. The WBC therefore connects the
accelerations and ground reaction forces produced by the DFCIP-MPC to the
complete multibody dynamics of TITA.\\
\
For TITA, the residual policy observes the robot state together with selected
outputs of the model-based controller. These include the accelerations and
ground reaction forces planned by MPC and the nominal torques computed by WBC.
The policy therefore has access not only to the measured motion but also to the
action intended by the nominal controller. Its outputs modify the desired
positions of the leg joints and the desired velocities of the wheels.
Joint-level feedback converts these offsets into residual torques, which are
added to the nominal WBC torques before the complete command is applied to the
actuators.\\
\
The second platform is Lite3, a quadruped robot controlled through an
SRBD-based MPC and a Whole-Body Controller. The MPC receives planar velocity
commands and predicts the motion of the robot base while optimizing the ground
reaction forces of the four feet over a finite horizon. A predefined gait
generator determines the stance and swing phases and constructs the
corresponding foot references. The WBC uses the first optimized set of contact
forces for stance control and tracks the planned foot trajectories during
swing, producing the nominal joint torques for all twelve actuators.\\
\
The Lite3 residual policy generates bounded offsets from a nominal joint
posture. A joint-space proportional--derivative controller converts these
offsets into residual torques, which are added to the output of the MPC--WBC
branch. The policy thus modifies the resulting whole-body behaviour without
replacing the planner or directly altering the MPC optimization. The Lite3
study also permits comparison with an end-to-end RL controller, making it
possible to examine the trade-off between a purely model-based solution, a
residual architecture, and a policy that learns locomotion without the same
online model-based prior.\\
\
Both residual policies are trained using Proximal Policy Optimization
\cite{ppo}, an on-policy policy-gradient algorithm that alternates
between collecting trajectories and performing multiple optimization epochs on
a clipped surrogate objective. Training is performed in parallel MuJoCo
environments, with the model-based controller active inside each environment.
The reward promotes tracking of the commanded linear and angular velocities
while penalizing unstable body motion, excessive actuator effort, rapid changes
in the learned action, and episode termination. During training, the critic may
use privileged simulation quantities that are not required by the deployed
actor. Observation noise and disturbances can also be introduced to reduce
dependence on ideal measurements and to test whether the learned correction
remains effective away from nominal conditions.\\
\
The objective of this thesis is to determine whether residual learning can
improve the tracking accuracy and robustness of an MPC-based locomotion
controller without discarding its physical structure. The investigation
considers individual and combined velocity commands, stability limits, control
effort, uneven terrain, disturbances, observation noise, and variations in
model parameters. The comparison is not limited to the final tracking error.
It also examines whether improvements are obtained through bounded and
sufficiently smooth corrections, whether the nominal behaviour is preserved
when no correction is needed, and whether the learned component introduces
systematic biases or undesirable interactions with the original controller.\\
\
The main contribution is the development and evaluation of two residual MPC
architectures for robots with different morphologies and contact behaviours.
The first combines a DFCIP-based MPC, a constrained WBC, and joint-level
residual control on a two-wheeled legged platform. The second applies the same
general principle to a quadruped using SRBD-MPC, gait scheduling, and
stance--swing whole-body control. In both cases, MPC remains active during
policy training and evaluation. This allows the policy to learn within the same
closed-loop architecture used at evaluation and provides a direct comparison
against the unmodified model-based controller. The study further analyses the
practical consequences of residual torque addition, including action
representation, controller update rates, baseline isolation, and the
interaction between the learned and nominal commands.\\
\
The remainder of the thesis develops this investigation from its theoretical
foundations to its experimental evaluation. The next chapter reviews
model-based, learning-based, and hybrid approaches to legged locomotion. The
robot models and the reduced-order representations used by the controllers are
then introduced, followed by the formulation of MPC and Whole-Body Control.
The reinforcement-learning problem, policy observations, actions, and training
procedure are subsequently described. The complete residual architectures for
TITA and Lite3 are then presented, including the exchange of information
between the model-based and learned components. Finally, the controllers are
evaluated in simulation, the results and limitations are discussed, and the
main conclusions and possible directions for future work are reported.
